% Unified Optimizer — Complete Mathematical Model
% Source of truth: src/optimizer/unified_optimizer.py
%                  src/optimizer/cost_matrix.py
%                  src/constraints/*.py
\documentclass[11pt,a4paper]{article}

\usepackage[margin=1in]{geometry}
\usepackage{amsmath,amssymb,amsfonts}
\usepackage{booktabs}
\usepackage{enumitem}
\usepackage{hyperref}
\usepackage{array}
\usepackage{longtable}

\hypersetup{colorlinks=true, linkcolor=blue, citecolor=blue, urlcolor=blue}

\title{Unified Vehicle Allocation and Charge Scheduling Optimizer\\[0.4em]
\large Complete Mathematical Model}
\author{allocation-v2 codebase}
\date{\today}

\newcommand{\R}{\mathbb{R}}
\newcommand{\Z}{\mathbb{Z}}
\newcommand{\1}{\mathbf{1}}

\begin{document}
\maketitle

\begin{abstract}
This document formalises the mixed-integer optimisation model implemented in the
\texttt{UnifiedOptimizer} (Hexaly). The system supports three operational modes:
\textbf{allocation-only}, \textbf{scheduling-only} (with pre-allocated routes), and
\textbf{integrated} (joint allocation and charging). Scheduling is solved primarily with
\textbf{interval variables} (one optional charging session per vehicle); a
\textbf{time-slot} formulation is used as fallback when interval support is unavailable.
Sequence scores for allocation are pre-computed by a constraint manager before the
solver is invoked.
\end{abstract}

\tableofcontents
\newpage

%==============================================================================
\section{Overview and Solver Architecture}
%==============================================================================

The unified optimizer solves a weighted-sum problem:
\begin{equation}
  \label{eq:integrated-objective}
  \max \;\; \alpha \cdot A(x) \;-\; \beta \cdot S(y),
\end{equation}
where $A(x)$ is the allocation score, $S(y)$ is the scheduling cost (including SOC
shortfall and timing penalties), $\alpha = \texttt{allocation\_score\_weight}$, and
$\beta = \texttt{scheduling\_cost\_weight}$.

\begin{itemize}[nosep]
  \item \textbf{Allocation-only:} maximise $A(x)$ subject to allocation constraints.
  \item \textbf{Scheduling-only:} minimise $S(y)$ subject to scheduling constraints
        (routes fixed via \texttt{fix\_allocation}).
  \item \textbf{Integrated:} maximise~\eqref{eq:integrated-objective} over $(x,y)$.
\end{itemize}

The implementation selects interval-based scheduling when Hexaly exposes
\texttt{interval\_var}; otherwise the time-slot model is used. Both variants share
pricing parameters ($\gamma_t$, $\sigma$, $\tau$) but differ in decision variables
and enforcement: the interval model adds an early-start penalty ($\mu$) and now
enforces per-slot site capacity via overlap-based average power; the time-slot model
uses per-slot SOC shortfall ($\lambda\sigma_{t,v}$) instead of interval-only final
shortfall, and reserves charger slots for all available assigned vehicles (plug-in
semantics).

%==============================================================================
\section{Sets, Indices, and Input Data}
%==============================================================================

\subsection{Index sets}

\begin{center}
\begin{tabular}{@{}ll@{}}
\toprule
Symbol & Description \\
\midrule
$V = \{1,\ldots,n_v\}$ & Vehicles (site fleet) \\
$R = \{1,\ldots,n_r\}$ & Routes in planning/allocation window \\
$S = \{1,\ldots,n_s\}$ & Feasible route \emph{sequences} (pre-generated) \\
$T = \{0,\ldots,n_t-1\}$ & Discrete 30-minute time slots \\
$K = \{0,\ldots,n_k-1\}$ & Charger power classes at the site \\
$P_v \subseteq R$ & Route checkpoints for vehicle $v$ (ordered by departure) \\
\bottomrule
\end{tabular}
\end{center}

Each sequence $s \in S$ is a tuple $(v(s), \mathcal{R}_s, c_s)$ where $v(s) \in V$ is
the assigned vehicle, $\mathcal{R}_s = (r_1,\ldots,r_{|\mathcal{R}_s|})$ is an ordered
route subsequence, and $c_s \in \R$ is the pre-computed sequence score.

\subsection{Static parameters}

\begin{center}
\begin{longtable}{@{}p{3.2cm}p{10cm}@{}}
\toprule
Symbol & Meaning \\
\midrule
\endhead
$W_{\mathrm{route}}$ & \texttt{route\_count\_weight} (default $10^2$): priority on covering routes \\
$\alpha$ & \texttt{allocation\_score\_weight} \\
$\beta$ & \texttt{scheduling\_cost\_weight} \\
$\lambda$ & \texttt{target\_soc\_shortfall\_penalty} (£ per kWh shortfall) \\
$\sigma$ & \texttt{synthetic\_time\_price\_factor} (early-charge bias) \\
$\tau$ & \texttt{triad\_penalty\_factor} (£ per kWh during TRIAD slots) \\
$\mu$ & \texttt{makespan\_penalty\_weight} (interval start-time penalty scale) \\
$\theta_v$ & Target SOC for vehicle $v$ (\%, converted to kWh) \\
$\Delta = 0.5$ & Slot duration (hours); each slot is 30 minutes \\
$H$ & Planning horizon in minutes ($H = 30\,n_t$) \\
\bottomrule
\end{longtable}
\end{center}

\subsection{Vehicle and site data}

For each vehicle $v \in V$:
\begin{itemize}[nosep]
  \item $B_v$ --- battery capacity (kWh)
  \item $s_v^0$ --- current SOC (kWh) at planning start
  \item $\bar{s}_v = (\theta_v/100)\,B_v$ --- target SOC (kWh)
  \item $\Delta s_v^{\mathrm{tgt}} = \max(0,\,\bar{s}_v - s_v^0)$ --- energy needed to reach target
  \item $\bar{p}_v$ --- AC charge rate limit (kW)
  \item $E_v^{\max} = B_v - s_v^0$ --- maximum deliverable energy
  \item $a_{v,t} \in \{0,1\}$ --- availability (1 = vehicle can charge in slot $t$)
\end{itemize}

For each slot $t \in T$:
\begin{itemize}[nosep]
  \item $\pi_t$ --- electricity price (£/kWh)
  \item $\delta_t \in \{0,1\}$ --- TRIAD indicator
  \item $d_t$ --- forecasted non-EV site demand (kW)
  \item $C = \texttt{site\_capacity\_kw}$ --- agreed site capacity (kW)
\end{itemize}

Slot unit cost (used in both scheduling formulations):
\begin{equation}
  \label{eq:slot-cost}
  \gamma_t = \pi_t \;+\; \sigma\,\frac{n_t - t}{n_t} \;+\; \tau\,\delta_t.
\end{equation}

For each route checkpoint $p \in P_v$ at slot index $t_p$:
\begin{itemize}[nosep]
  \item $e_p$ --- cumulative route energy requirement at departure (kWh)
  \item Required delivered energy before departure:
        $\rho_{v,p} = \max(0,\, e_p - s_v^0)$
\end{itemize}

For each charger class $k \in K$:
\begin{itemize}[nosep]
  \item $N_k$ --- number of physical chargers in class $k$
  \item $\bar{P}_k$ --- maximum power (kW) of class $k$
\end{itemize}

%==============================================================================
\section{Pre-processing: Sequence Generation and Scores}
%==============================================================================

Before the Hexaly model is built, \texttt{CostMatrixBuilder} enumerates feasible
route sequences per vehicle. For vehicle $v$, all subsets of routes of length
$1,\ldots,L_v$ are considered (combinatorial generation with overlap pruning).

\subsection{Hard feasibility (sequence filter)}

A sequence $\mathcal{R} = (r_1,\ldots,r_m)$ for vehicle $v$ is \textbf{discarded} if any
\emph{hard} MAF constraint evaluates to a negative penalty. Hard constraints include:
\begin{itemize}[nosep]
  \item \textbf{Energy feasibility} --- sufficient battery energy for the route sequence
  \item \textbf{Strict turnaround} --- consecutive routes do not overlap in time
        (modulo configured turnaround minutes)
  \item \textbf{Strict shift hours} --- routes fall within driver shift windows
\end{itemize}

Non-overlapping route pairs are also pre-filtered during sequence generation using
turnaround-aware overlap checks.

\subsection{Soft sequence score}

For each feasible sequence $\mathcal{R}$ assigned to vehicle $v$, the score is:
\begin{equation}
  \label{eq:sequence-score}
  c_{v,\mathcal{R}} = \sum_{j \in \mathcal{J}} \mathrm{cost}_j(v, \mathcal{R}),
\end{equation}
where $\mathcal{J}$ is the set of enabled soft constraints (e.g.\ preferred turnaround,
charger preference bonuses/penalties). Hard-constraint violations yield
$c = -\infty$ (sequence excluded).

Each surviving sequence becomes one index $s \in S$ with score $c_s = c_{v(s),\mathcal{R}_s}$.

%==============================================================================
\section{Allocation Sub-model}
%==============================================================================

\subsection{Decision variables}

\begin{align}
  x_s &\in \{0,1\} && \text{1 if sequence $s$ is selected} \\
  y_r &\in \{0,1\} && \text{1 if route $r$ is covered (auxiliary)}
\end{align}

\subsection{Objective (allocation-only)}

\begin{equation}
  \label{eq:alloc-obj}
  \max \;\; W_{\mathrm{route}} \sum_{r \in R} y_r \;+\; \sum_{s \in S} c_s\, x_s.
\end{equation}

\subsection{Constraints}

\paragraph{At most one sequence per vehicle.}
Let $S_v = \{s \in S : v(s) = v\}$.
\begin{equation}
  \sum_{s \in S_v} x_s \;\le\; 1 \qquad \forall v \in V.
\end{equation}

\paragraph{Route coverage linking.}
For each route $r \in R$, let $S(r) = \{s \in S : r \in \mathcal{R}_s\}$.
\begin{align}
  \sum_{s \in S(r)} x_s &\;\le\; 1 && \forall r \in R, \\
  y_r &\;\le\; \sum_{s \in S(r)} x_s && \forall r \in R, \\
  \sum_{s \in S(r)} x_s &\;\le\; |S(r)|\, y_r && \forall r \in R.
\end{align}

The first inequality ensures no route is allocated twice; the linking constraints
define $y_r = \1[\sum_{s \in S(r)} x_s \ge 1]$.

%==============================================================================
\section{Scheduling Sub-model}
%==============================================================================

Two equivalent scheduling economics are implemented; the solver uses one formulation
depending on Hexaly capabilities.

%------------------------------------------------------------------------------
\subsection{Formulation A: Interval Model (Primary)}
%------------------------------------------------------------------------------

Each vehicle $v$ has one optional charging session represented by an interval variable.

\subsubsection{Decision variables}

\begin{align}
  I_v &= [\mathrm{start}_v,\;\mathrm{end}_v) && \text{charging session interval (minutes)} \\
  \ell_v &= \mathrm{end}_v - \mathrm{start}_v && \text{session duration (minutes)} \\
  E_v &\in [0,\, E_v^{\max}] && \text{total energy delivered (kWh)} \\
  k_v &\in K && \text{charger power class (if allocation enabled)}
\end{align}

Optional sessions satisfy $0 \le \ell_v \le \ell_v^{\max}$ where
$\ell_v^{\max} = \lfloor 60\, E_v^{\max} / \bar{p}_v \rfloor$.

Energy-duration coupling at effective charge rate $\min(\bar{p}_v, \bar{P}_{k_v})$:
\begin{equation}
  \label{eq:energy-duration}
  E_v = \frac{\min(\bar{p}_v, \bar{P}_{k_v})}{60}\,\ell_v.
\end{equation}
When charger allocation is disabled, $\bar{P}_{k_v}$ is omitted and $\bar{p}_v$ alone
is used.

Fixed slot intervals for costing: for each $t \in T$,
$J_t = [30t,\, 30(t+1))$ (minutes).

Overlap energy in slot $t$:
\begin{equation}
  \label{eq:overlap-energy}
  e_{v,t} = \frac{\min(\bar{p}_v, \bar{P}_{k_v})}{60}\,\mathrm{overlap\_length}(I_v, J_t).
\end{equation}

\subsubsection{Scheduling objective (scheduling-only)}

Minimise:
\begin{equation}
  \label{eq:sched-interval-obj}
  \underbrace{\sum_{v \in V}\sum_{t \in T} \gamma_t\, e_{v,t}}_{\text{slot-aware charging cost}}
  \;+\;
  \underbrace{\frac{\mu}{H}\sum_{v \in V} \mathrm{start}_v}_{\text{early-start penalty}},
\end{equation}
The interval path uses hard target SOC (eq:hard-target) rather than a final
shortfall term; $\sigma_v^{\mathrm{fin}}$ is omitted because it would be zero at
optimality whenever the hard constraint is active.
\begin{equation}
  \label{eq:hard-target}
  E_v \;\ge\; \Delta s_v^{\mathrm{tgt}} \qquad \forall v \text{ with } \Delta s_v^{\mathrm{tgt}} > 0.
\end{equation}

\subsubsection{Scheduling constraints (interval)}

\paragraph{Route energy (hard).}
For each checkpoint $p \in P_v$ with departure at minute $m_p$ and requirement $\rho_{v,p} > 0$:
\begin{align}
  \mathrm{end}_v &\;\le\; m_p, \\
  E_v &\;\ge\; \rho_{v,p}.
\end{align}
\emph{Expressiveness note:} with a single optional charging interval per vehicle,
all $\rho_{v,p}$ must be satisfied before the \textbf{earliest} relevant departure;
inter-leg charging between multiple route legs is not modeled.

\paragraph{Site capacity (interval).}
Per slot $t$ with available capacity $C - d_t$:
\begin{equation}
  \label{eq:site-capacity-interval}
  \sum_{v \in V} \frac{\min(\bar{p}_v, \bar{P}_{k_v})}{30}\,
  \mathrm{overlap\_length}(I_v, J_t)
  \;\le\; \max(0,\, C - d_t).
\end{equation}
Average slot power is used for consistency with eq:overlap-energy. A startup warning
is logged when $\sum_k N_k \bar{P}_k > C$ (nameplate exceeds contract).

\paragraph{Availability window.}
Let $t_v^{\mathrm{first}}$ and $t_v^{\mathrm{last}}$ be the first and last available slot
indices for vehicle $v$. The session must lie within the contiguous available span:
\begin{align}
  \mathrm{start}_v &\;\ge\; 30\, t_v^{\mathrm{first}}, \\
  \mathrm{end}_v &\;\le\; 30\,(t_v^{\mathrm{last}} + 1).
\end{align}

\paragraph{Integrated mode: route execution intervals.}
When allocation and scheduling are solved jointly, each selected sequence $s$ induces
route execution intervals $R_{s,j}$ for $j = 1,\ldots,|\mathcal{R}_s|$. For selected
sequences:
\begin{align}
  \mathrm{present}(R_{s,j}) &\;\Leftrightarrow\; x_s = 1, \\
  \mathrm{start}(R_{s,j+1}) &\;\ge\; \mathrm{end}(R_{s,j}), \\
  \mathrm{no\_overlap}(I_{v(s)}, R_{s,j}) &\quad \forall j.
\end{align}

\paragraph{Charger class assignment (optional).}
When \texttt{enable\_charger\_allocation} is true:
\begin{itemize}[nosep]
  \item Each vehicle selects exactly one class: $k_v \in K$.
  \item If already connected to a charger, $k_v$ is fixed to that charger's class.
  \item Nighttime continuity: vehicles that charged overnight on class $k$ must use
        the same class if charging again during night hours ($19{:}00$--$07{:}00$).
  \item \textbf{Cumulative capacity} via pulse functions: for each class $k$,
        \begin{equation}
          \sum_{v \in V} \mathrm{pulse}\bigl(I_v,\; \1[k_v = k]\bigr) \;\le\; N_k.
        \end{equation}
        The pulse is active while vehicle $v$ is charging and assigned to class $k$.
  \item Effective rate $\min(\bar{p}_v, \bar{P}_{k_v})$ is enforced in
        eq:energy-duration and eq:overlap-energy (not a separate conditional block).
\end{itemize}

%------------------------------------------------------------------------------
\subsection{Formulation B: Time-Slot Model (Fallback)}
%------------------------------------------------------------------------------

\subsubsection{Decision variables}

\begin{align}
  p_{t,v} &\in [0,\,\bar{p}_v] && \text{charging power in slot $t$ (kW)} \\
  u_{t,v} &\in [0,\, E_v^{\max}] && \text{cumulative energy delivered by end of slot $t$ (kWh)} \\
  z_{v,k} &\in \{0,1\} && \text{1 if vehicle $v$ assigned to charger class $k$}
\end{align}

\subsubsection{Cumulative energy dynamics}

\begin{equation}
  u_{t,v} =
  \begin{cases}
    \Delta\, p_{0,v} & t = 0, \\[4pt]
    u_{t-1,v} + \Delta\, p_{t,v} & t > 0.
  \end{cases}
\end{equation}

Energy delivered in slot $t$: $e_{t,v} = \Delta\, p_{t,v}$.

\subsubsection{Scheduling objective (time-slot)}

Minimise:
\begin{equation}
  \label{eq:sched-timeslot-obj}
  \sum_{t \in T}\sum_{v \in V} \gamma_t\, e_{t,v}
  \;+\;
  \frac{\lambda}{n_t}\sum_{t \in T}\sum_{v \in V} \sigma_{t,v}
  \;+\;
  \text{(integrated: } \alpha,\beta \text{ weighting as in \eqref{eq:integrated-objective})},
\end{equation}
where per-slot shortfall variables satisfy:
\begin{align}
  \sigma_{t,v} &\;\ge\; 0 && \forall t,v, \\
  \sigma_{t,v} &\;\ge\; \Delta s_v^{\mathrm{tgt}} - u_{t,v} && \forall t,v \text{ with } \Delta s_v^{\mathrm{tgt}} > 0.
\end{align}

Hard target SOC at horizon end:
\begin{equation}
  u_{n_t-1,v} \;\ge\; \Delta s_v^{\mathrm{tgt}} \qquad \forall v \text{ with } \Delta s_v^{\mathrm{tgt}} > 0.
\end{equation}

\subsubsection{Scheduling constraints (time-slot)}

\paragraph{Route energy (hard).}
For checkpoint $p$ at slot $t_p$ with $\rho_{v,p} > 0$:
\begin{equation}
  u_{t_p - 1,\, v} \;\ge\; \rho_{v,p}.
\end{equation}

\paragraph{Site capacity.}
\begin{equation}
  \sum_{v \in V} p_{t,v} \;\le\; \max(0,\, C - d_t) \qquad \forall t \in T.
\end{equation}

\paragraph{Battery ceiling.}
\begin{equation}
  u_{t,v} \;\le\; E_v^{\max} \qquad \forall t \in T,\; \forall v \in V.
\end{equation}

\paragraph{Charge rate.}
\begin{equation}
  p_{t,v} \;\le\; \bar{p}_v \qquad \forall t,v.
\end{equation}

\paragraph{Availability.}
\begin{equation}
  p_{t,v} = 0 \qquad \forall t \text{ with } a_{v,t} = 0.
\end{equation}

\paragraph{Charger class assignment (optional).}
When enabled:
\begin{align}
  \sum_{k \in K} z_{v,k} &= 1 && \forall v, \\
  p_{t,v} &\;\le\; \sum_{k \in K} \min(\bar{p}_v, \bar{P}_k)\, z_{v,k} && \forall t,v, \\
  \sum_{v:\, a_{v,t}=1} z_{v,k} &\;\le\; N_k && \forall t,k.
\end{align}
The charger-count constraint uses \textbf{plug-in reservation}: any available vehicle
assigned to class $k$ counts against $N_k$, even when $p_{t,v} = 0$ (idle but connected).
Fixed-class and nighttime-continuity rules mirror the interval formulation.

%==============================================================================
\section{Integrated Model}
%==============================================================================

The integrated model combines Section~\ref{eq:alloc-obj} allocation variables and
constraints with Section~5 scheduling variables and constraints. The objective is
\begin{equation}
  \max \;\; \alpha \left( W_{\mathrm{route}} \sum_{r \in R} y_r + \sum_{s \in S} c_s x_s \right)
  - \beta \left( \text{scheduling cost terms from \eqref{eq:sched-interval-obj} or \eqref{eq:sched-timeslot-obj}} \right).
\end{equation}

Route execution intervals are linked to sequence selection via
$\mathrm{present}(R_{s,j}) \Leftrightarrow x_s = 1$, and charging sessions must not
overlap allocated route intervals for the same vehicle.

%==============================================================================
\section{Objective Component Summary}
%==============================================================================

\begin{center}
\begin{tabular}{@{}llll@{}}
\toprule
Component & Interval model & Time-slot model & Role \\
\midrule
Electricity cost & $\sum_{v,t} \gamma_t e_{v,t}$ & $\sum_{v,t} \gamma_t \Delta p_{t,v}$ & TOU + TRIAD pricing \\
Early-charge bias & $\sigma$ in $\gamma_t$; $\mu$ start penalty & $\sigma$ in $\gamma_t$ only & Prefer earlier charging \\
Site capacity & eq:site-capacity-interval & eq:site-capacity (time-slot) & Cap aggregate draw \\
SOC shortfall (soft) & --- (hard target only) & $\frac{\lambda}{n_t}\sum_{t,v}\sigma_{t,v}$ & Penalise below-target SOC \\
SOC target (hard) & $E_v \ge \Delta s_v^{\mathrm{tgt}}$ & $u_{n_t-1,v} \ge \Delta s_v^{\mathrm{tgt}}$ & Require target when feasible \\
Route energy (hard) & $E_v \ge \rho_{v,p}$, $\mathrm{end}_v \le m_p$ & $u_{t_p-1,v} \ge \rho_{v,p}$ & Feasible departures \\
Allocation score & $\sum_s c_s x_s + W_{\mathrm{route}}\sum_r y_r$ & (same) & Route assignment quality \\
\bottomrule
\end{tabular}
\end{center}

%==============================================================================
\section{Default Parameter Values}
%==============================================================================

\begin{center}
\begin{tabular}{@{}lll@{}}
\toprule
Parameter & Config key & Default \\
\midrule
$W_{\mathrm{route}}$ & \texttt{route\_count\_weight} & $10^2$ \\
$\alpha$ & \texttt{allocation\_score\_weight} & $1.0$ \\
$\beta$ & \texttt{scheduling\_cost\_weight} & $1.0$ \\
$\lambda$ & \texttt{target\_soc\_shortfall\_penalty} & $2.0$ \\
$\sigma$ & \texttt{synthetic\_time\_price\_factor} & $0.05$ \\
$\tau$ & \texttt{triad\_penalty\_factor} & $100.0$ \\
$\mu$ & \texttt{makespan\_penalty\_weight} & $0.5$ \\
$\theta_v$ & \texttt{target\_soc\_percent} & $75\%$ config default; $85\%$ via MAF (\texttt{DEFAULT\_TARGET\_SOC\_PERCENT} in \texttt{src/config.py}, resolved in \texttt{get\_scheduler\_params\_from\_site\_config}) \\
\bottomrule
\end{tabular}
\end{center}

%==============================================================================
\section{Implementation Notes}
%==============================================================================

\begin{enumerate}[label=\arabic*.]
  \item \textbf{Solver:} Hexaly Optimizer (\texttt{hexaly.optimizer}). Time limits:
        30\,s (allocation), 300\,s (scheduling), 330\,s (integrated).
  \item \textbf{Greedy fallback:} When Hexaly is inactive or the solve throws, a greedy
        heuristic selects cheapest available slots (by $\gamma_t$) to meet energy needs.
  \item \textbf{Sequence explosion:} $n_s$ grows combinatorially in $|R|$ and
        \texttt{max\_routes\_per\_vehicle}; infeasible sequences are filtered before
        the solver sees them.
  \item \textbf{Interval availability simplification:} The interval model uses the
        first-to-last available slot span (contiguous envelope), not individual gaps;
        the time-slot model respects per-slot availability exactly. Post-solve logging
        flags sessions that overlap unavailable slots inside the envelope.
  \item \textbf{Single charging session:} One optional interval per vehicle cannot
        split energy across inter-leg windows; see route-energy expressiveness note.
  \item \textbf{Site capacity validation:} If $\sum_k N_k \bar{P}_k > C$, a warning
        is logged at model build; per-slot constraints still enforce $C - d_t$.
  \item \textbf{API overrides:} All weights and $\theta_v$ can be overridden per request
        via \texttt{/optimize/unified}.
\end{enumerate}

%==============================================================================
\section{Notation Quick Reference}
%==============================================================================

\begin{center}
\begin{tabular}{@{}ll@{}}
\toprule
Symbol & Definition \\
\midrule
$x_s$ & Sequence $s$ selected \\
$y_r$ & Route $r$ covered \\
$p_{t,v}$ & Charging power (kW) in slot $t$ \\
$u_{t,v}$ & Cumulative energy delivered (kWh) by end of slot $t$ \\
$E_v$ & Total energy delivered in interval session (kWh) \\
$I_v$ & Charging session interval \\
$\gamma_t$ & Slot unit cost \eqref{eq:slot-cost} \\
$\Delta s_v^{\mathrm{tgt}}$ & Energy gap to target SOC \\
$\rho_{v,p}$ & Energy gap to route checkpoint $p$ \\
$c_s$ & Pre-computed sequence score \eqref{eq:sequence-score} \\
\bottomrule
\end{tabular}
\end{center}

\end{document}
